\par\bigskip\Needspace{7\baselineskip}
\pdfbookmark[1]{Theory of Biquadratic Residues. III}{nachlass-biquadratic-III}
\begin{center}[III.]\par\medskip\textsc{Theory of Biquadratic Residues.}\end{center}
\BookRule{.08\textwidth}
\renewcommand{\Article}[1]{\par\medskip\Needspace{4\baselineskip}\pdfbookmark[2]{Art. #1}{nachlass-biquadratic-III-art-#1}\centerline{\textbf{#1.}}\nobreak\smallskip}
\Article{1}
The least residues for the modulus $m=a+bi$ are the complex integers $\mu=\alpha+\beta i$ for which $\frac{\mu}{m}=x+yi$ has $x$ and $y$ positive and less than $1$. The residue $0$ is included as well\begingroup\renewcommand{\thefootnote}{*)}\footnote{and, if $a$ and $b$ happen to have a common divisor $e$, the numbers $\frac{m}{e},\frac{2m}{e},\frac{3m}{e}\dotstwo\frac{(e-1)m}{e}$. For now, however, we shall leave this case out of the investigation.}\endgroup. Their number is $=aa+bb$.

\Article{2}
If $aa+bb$ is odd, $aa+bb$ is of the form $4n+1$. On excluding the least residue $0$, the others fall into four classes. In the first class $f$ we place those for which $x$ and $y$ are less than $\frac12$,
\SourceLeaf{342}
\[
\begin{array}{lcll}
\text{the second}&f'&\text{where }x>\tfrac12,&y<\tfrac12\\
\text{third}&f''&x>\tfrac12,&y>\tfrac12\\
\text{fourth}&f'''&x<\tfrac12,&y>\tfrac12
\end{array}
\]
All the residues are obtained
\[
\begin{array}{rcl}
f'&\text{from}&if+m\\
f''&\text{from}&-f+(1+i)m\\
f'''&\text{from}&-if+im
\end{array}
\]

\Article{3}
Let $M$ be another number with no factor in common with $m$. Then
\[
M^{aa+bb-1}\equiv1\ (\mathrm{mod.}\ m)
\]
holds; consequently $M^{\frac14(aa+bb-1)}$ is either $\equiv1$, or $\equiv i$, or $\equiv-1$, or $\equiv-i$, that is, $\equiv i^{\varepsilon}$, where $\varepsilon$ represents one of the four numbers $0,1,2,3$. In the first case $M$ is a biquadratic residue of $m$, and hence also a quadratic residue. In the third case $M$ is a quadratic residue but not a biquadratic residue; in the second and fourth cases it is both a quadratic and a biquadratic nonresidue. We call this $\varepsilon$, on which the biquadratic character of the number $M$ with respect to the modulus $m$ depends, the decident of $M$ for the modulus $m$. Induction gives the following beautiful theorem. “If $M$ and $m$ are odd primes of the form $1+(2+2i)\mu$, where $\mu$ is an integer, the difference between the two decidents, that of $M$ for the modulus $m$ and that of $m$ for the modulus $M$, is either $=0$ or $=2$; the former when at least one of the numbers $m,M$ is of the form $1+4N$, the latter when both are of the form $1+2i+4N$.” This reciprocity theorem is analogous to the one for quadratic residues among real numbers alone.

\Article{4}
Multiply all the numbers $f$ by $M$ and find their least residues for the modulus $m$. If among them there are
\[
\begin{array}{cll}
\alpha&\text{belonging to }f&\text{}\\
\beta&f'&\\
\gamma&f''&\\
\delta&f'''&
\end{array}
\]
then $\varepsilon\equiv\beta+2\gamma+3\delta\ (\mathrm{mod.}\ 4)$.

\SourceLeaf{343}
Proof. Let the collection of those numbers in $f$ whose products with $M$ give residues belonging to $f$ be $g$; let the collection of those whose products give residues from $f'$ be $g'$, and similarly $g''\pdot g'''$. Then the least residues of
\[
-ig'M,\quad -g''M,\quad ig'''M
\]
are all contained in $f$ and are distinct both from one another and from the least residues of the products $gM$. Consequently the product of all
\[
gM,\quad -ig'M,\quad -g''M,\quad +ig'''M
\]
is congruent to the product of all $f$, and hence also to the product of all $g,g',g'',g'''$. But that product equals the product of all $g,g',g'',g'''$ multiplied by
\[
M^{\alpha}\pdot(-iM)^{\beta}\pdot(-M)^{\gamma}\pdot(iM)^{\delta}
\]
so this last product is $\equiv1$
\[
\begin{aligned}
\text{therefore}\qquad&M^{\alpha+\beta+\gamma+\delta}(-i)^{\beta}(-1)^{\gamma}i^{\delta}\equiv1\\
\text{or}\qquad&M^{\alpha+\beta+\gamma+\delta}\equiv i^{\beta}(-1)^{\gamma}(-i)^{\delta}\equiv i^{\beta+2\gamma+3\delta}
\end{aligned}
\]
from which the theorem follows directly.

\Article{5}
It is easy to decide whether the least residue of a number $N$ for the modulus $m$ belongs to the class $f,f',f''$ or $f'''$. If $\omega$ is the integer contained in $\frac{N}{m}$, that residue is $=N-\omega m$ and therefore belongs to $f,f',f''$, according as, on putting
\[
\frac{N}{m}-\omega=x+iy
\]
we have
\[
\begin{array}{ll}
x<\tfrac12,&y<\tfrac12\\
x>\tfrac12,&y<\tfrac12\\
x>\tfrac12,&y>\tfrac12\\
x<\tfrac12,&y>\tfrac12
\end{array}
\]
In these $4$ cases the integer contained in $\frac{2N}{m}$ is, in the same order, as follows:
\SourceLeaf{344}
\[
\begin{array}{l}
2\omega\\
2\omega+1\\
2\omega+1+i\\
2\omega+i
\end{array}
\]
It follows that the least residue of $N$ for the modulus $m$ belongs to $f,f',f'',f'''$ according as, on writing the integer contained in $\frac{2N}{m}$ as $=\xi+\eta i$, we have
\[
\begin{array}{ll}
\xi\ \text{even}&\eta\ \text{even}\\
\xi\ \text{odd}&\eta\ \text{even}\\
\xi\ \text{odd}&\eta\ \text{odd}\\
\xi\ \text{even}&\eta\ \text{odd}
\end{array}
\]

\Article{6}
The decident of $M$ for the modulus $m$ is therefore found as follows. Find the integers contained in each of the individual $\frac{2fM}{m}$. Writing these generally as $x+yi$, disregard entirely those for which $x$ and $y$ are both even; count one for each for which $x$ is odd and $y$ even, take two for each for which $x$ and $y$ are both odd, and three for each for which $x$ is even and $y$ odd. Take the least residue modulo $4$ of the sum of all these numbers; it is the required decident. We express this as
\[
\mathrm{Dec.}\frac{M}{m}=\Sigma n
\]
where $[\frac{2fM}{m}]=x+yi$, and we put $n=0$ when $x$ is even and $y$ even
\[
\begin{array}{cll}
1&x\ \text{odd}&y\ \text{even}\\
2&x\ \text{odd}&y\ \text{odd}\\
3&x\ \text{even}&y\ \text{odd}
\end{array}
\]
For brevity we shall denote $n$ by the characteristic $\Theta$, $n=\Theta\frac{2fM}{m}$\begingroup\renewcommand{\thefootnote}{*)}\footnote{To decide which class $M$ belongs to with respect to $m$, choose those values of $k$ (among the numbers $1,2,3\ldots p-1$) for which $[\frac{2km'M}{p}]$ is even, and add $-\Sigma[\frac{2km'}{p}]^2$. If $k$ is taken only up to $\frac12p$, the sum to be taken is
\[
-\Sigma\{[\frac{2km'M}{p}]^2+[\frac{2km'}{p}]^2\}
\]
for those values of $[\frac{2km'M}{p}]$ that are divisible by $1+i$.}\endgroup.
